%!TEX program = xelatex
\documentclass[cn,blue,14pt,normal]{elegantnote}

\usepackage{chemformula} %化学公式宏包
\usepackage{siunitx} %数字，单位宏包
\sisetup{
	separate-uncertainty = true,
	inter-unit-product = \ensuremath{{}\cdot{}}
} %单位间加小圆点
\usepackage{tikz} %Tikz绘图包
\usetikzlibrary{cd}
\usetikzlibrary{chains}
\usetikzlibrary{arrows}
\usetikzlibrary{arrows.meta} %画箭头
\usetikzlibrary{commutative-diagrams} %交换图
\usetikzlibrary{positioning,automata}
\usepackage[all]{xy}
\usepackage{booktabs} %三线表
\usepackage{multirow} %合并单元格
\usepackage{makecell} %表格内强制换行
\usepackage{array}
\usepackage{booktabs} %控制表格行高
\usepackage{colortbl}  %彩色表格需要加载的宏包
\usepackage{amsmath}
\usepackage{unicode-math}
\usepackage{fontspec} %字体

\setmainfont{Times New Roman}
\setmathfont{XITS Math}

%微分符号
\makeatletter
\newcommand*{\dif}{\mathop{}\!\mathrm{d}}
\makeatother

%直立积分
\DeclareSymbolFont{EulerExtension}{U}{euex}{m}{n}
\DeclareMathSymbol{\euintop}{\mathop} {EulerExtension}{"52}
\DeclareMathSymbol{\euointop}{\mathop} {EulerExtension}{"48}
\let\intop\euintop
\let\ointop\euointop

\begin{document}

习题2.9：在一恒容反应器中进行下列液相反应 ：
        \[
            \ch{A + B -> R} \qquad r_\mathrm{R} = 1.6 c_\mathrm{A}
        \]
        \[
            \ch{2 A -> D}  \qquad r_\mathrm{D} = 8.2 c_\mathrm{A}^2
        \]
        
        式中$r_\mathrm{R}$、$r_\mathrm{D}$分别表示产物 R 及 D 的生成速率，其单位均为 \si{kmol.m^{-3}.h^{-1}}， 反应用的原料为 A 与 B 的混合物 ， 其中 A 的浓度为 \SI{2}{kmol.m^{-3}} ， 试计算 A 的转化率达 95\%时所需要的反应时间 。

解：
\[
    - \mathcal{R}_\mathrm{A} = r_\mathrm{R} + 2 r_\mathrm{D} = 1.6 c_\mathrm{A} + 16.4 c_\mathrm{A}^2
\]

\[
    -\frac{\dif c_\mathrm{A}}{\dif t} = 1.6 c_\mathrm{A} + 16.4 c_\mathrm{A}^2
\]

移项：
\[
    -\frac{1}{c_\mathrm{A} + 10.25 c_\mathrm{A}^2}\dif c_\mathrm{A} = 1.6 \dif t
\]

积分：
\[
    -\int_{c_\mathrm{A0}}^{c_\mathrm{A}}\frac{1}{c_\mathrm{A} + 10.25 c_\mathrm{A}^2}\dif c_\mathrm{A} = \int_{0}^{t}1.6 t
\]

\[
    \int_{c_\mathrm{A}}^{c_\mathrm{A0}}\frac{1}{c_\mathrm{A} + 10.25 c_\mathrm{A}^2}\dif c_\mathrm{A} = \int_{0}^{t}1.6 t
\]

\[
    \int_{0.1}^{2}\frac{1}{c_\mathrm{A} + 10.25 c_\mathrm{A}^2}\dif c_\mathrm{A} = 1.6 t
\]

\[
    \int_{0.1}^{2}\frac{1}{c_\mathrm{A}(1 + 10.25 c_\mathrm{A})}\dif c_\mathrm{A} = 1.6 t
\]

\[
    \int_{0.1}^{2}\left(\frac{1}{c_\mathrm{A}} - \frac{10.25}{1 + 10.25 c_\mathrm{A}}\right)\dif c_\mathrm{A} = 1.6 t
\]

\[
    \int_{0.1}^{2}\left(\frac{1}{c_\mathrm{A}}\dif c_\mathrm{A} - \frac{10.25}{1 + 10.25 c_\mathrm{A}}\dif c_\mathrm{A}\right) = 1.6 t
\]

\[
    \int_{0.1}^{2}\left[\ln c_\mathrm{A} - \frac{1}{1 + 10.25 c_\mathrm{A}}\dif (1 + 10.25 c_\mathrm{A})\right] = 1.6 t
\]

\[
    \int_{0.1}^{2}\left(\ln c_\mathrm{A} - \ln (1 + 10.25 c_\mathrm{A})\right) = 1.6 t
\]

\[
    [\ln(2) - \ln(1+10.25\times 2)] - [\ln(0.1) - \ln(1+10.25\times 0.1)] = 1.6t
\]

\[
    t = 0.3958
\]









\end{document}
